% =====================================================================
%  CT Dose AI — Medical & Physics Cheatsheet
%  Compile with:  pdflatex CT_Dose_AI_Cheatsheet.tex   (run twice)
% =====================================================================
\documentclass[9pt,a4paper]{extarticle}

\usepackage[LGR,T1]{fontenc}
\usepackage[utf8]{inputenc}
\usepackage{helvet}\renewcommand{\familydefault}{\sfdefault}
\usepackage{microtype,amsmath,amssymb}
\usepackage{textalpha}   % lets UTF-8 Greek (ΥΤ ΣΤΕΦΑΝΙΟΓΡΑΦΙΑ ...) compile with pdflatex
\usepackage[a4paper,top=13mm,bottom=14mm,left=12mm,right=12mm]{geometry}
\usepackage{multicol,fancyhdr,titlesec,enumitem,parskip}
\usepackage{booktabs,tabularx,array,colortbl}
\usepackage{xcolor,tcolorbox}
\usepackage[hidelinks]{hyperref}
\tcbuselibrary{skins,breakable}

% ---------- palette ----------
\definecolor{ctblue}{HTML}{2A78D6}
\definecolor{ctbluedark}{HTML}{1A4F94}
\definecolor{ctaqua}{HTML}{1BAF7A}
\definecolor{ctaquadark}{HTML}{127A55}
\definecolor{ctorange}{HTML}{EB6834}
\definecolor{ctink}{HTML}{0B0B0B}
\definecolor{ctink2}{HTML}{52514E}
\definecolor{ctline}{HTML}{E2E2DE}
\definecolor{ctsurf}{HTML}{F7F7F5}
\definecolor{ctgreenbg}{HTML}{E8F6F0}

\newcommand{\secbar}[1]{\colorbox{ctbluedark}{\parbox{\dimexpr\columnwidth-2\fboxsep\relax}{%
  \color{white}\strut #1}}}
\titleformat{\section}{\normalsize\bfseries}{}{0pt}{\secbar}
\titlespacing*{\section}{0pt}{10pt}{5pt}
\titleformat{\subsection}{\small\bfseries\color{ctbluedark}}{}{0pt}{}
\titlespacing*{\subsection}{0pt}{7pt}{2pt}
\setlist[itemize]{leftmargin=10pt,itemsep=1pt,topsep=2pt,parsep=0pt}

\pagestyle{fancy}\fancyhf{}
\renewcommand{\headrulewidth}{0pt}\renewcommand{\footrulewidth}{0.4pt}
\fancyfoot[L]{\scriptsize\color{ctink2}CT Dose AI \textperiodcentered\ medical \& physics cheatsheet}
\fancyfoot[R]{\scriptsize\color{ctink2}\thepage}

\newcommand{\ctdi}{\ensuremath{\mathrm{CTDI}_{\mathrm{vol}}}}
\newcommand{\mgycm}{mGy$\cdot$cm}
% term entry:  \T{term}{definition}
\newcommand{\T}[2]{\par\smallskip\noindent{\color{ctbluedark}\bfseries #1} \ --- \ #2\par}
\newcommand{\Tg}[3]{\par\smallskip\noindent{\color{ctbluedark}\bfseries #1} {\scriptsize\color{ctink2}(#2)} \ --- \ #3\par}
\newcommand{\note}[1]{{\scriptsize\color{ctink2}#1\par}}

\newtcolorbox{keybox}{enhanced,breakable,colback=ctgreenbg,colframe=ctaqua,boxrule=0pt,leftrule=2pt,
  arc=1.5pt,left=6pt,right=6pt,top=4pt,bottom=4pt,fontupper=\small}
\newtcolorbox{warnbox}{enhanced,breakable,colback=ctsurf,colframe=ctorange,boxrule=0pt,leftrule=2pt,
  arc=1.5pt,left=6pt,right=6pt,top=4pt,bottom=4pt,fontupper=\small}
\newtcolorbox{eqbox}{enhanced,breakable,colback=ctsurf,colframe=ctline,boxrule=0.4pt,arc=2pt,
  left=6pt,right=6pt,top=4pt,bottom=4pt,fontupper=\small}

\newcolumntype{L}{>{\raggedright\arraybackslash}X}
\newcolumntype{R}{>{\raggedleft\arraybackslash}X}
\renewcommand{\arraystretch}{1.15}
\newcommand{\hdr}{\rowcolor{ctsurf}}

\begin{document}

% ---------------- title band ----------------
\begin{tcolorbox}[colback=ctbluedark,colframe=ctbluedark,arc=3pt,left=10pt,right=10pt,top=7pt,bottom=7pt]
{\color{white}\Large\bfseries CT Dose AI --- medical \& physics cheatsheet}\\[2pt]
{\color{white!85}\small Every clinical term used in the project, what it means, why it matters, and the typical
numbers. Written for an engineer who needs to talk to radiologists and medical physicists.}
\end{tcolorbox}

\begin{multicols}{2}

% =====================================================================
\section{1 \ \ The examination}

\Tg{CT}{computed tomography}{An X-ray tube and a detector array rotate around the patient. Many projections are
mathematically reconstructed into cross-sectional images. Unlike a plain X-ray it gives 3-D information --- at a
higher radiation dose.}

\Tg{CCTA}{coronary CT angiography; Greek: ΥΤ ΣΤΕΦΑΝΙΟΓΡΑΦΙΑ}{A CT of the heart's own arteries (the coronary
arteries) after intravenous iodine contrast. It answers: are the coronary arteries narrowed or blocked? It is the
examination in our dataset. A non-invasive alternative to catheter angiography.}

\Tg{Contrast agent}{iodinated contrast}{An iodine-containing liquid injected into a vein. Iodine absorbs X-rays
strongly, so blood vessels light up (become bright) against surrounding tissue. Without it, a coronary artery is
invisible. Iodine is the reason low kV works so well --- see §4.}

\Tg{Calcium score}{CAC, Agatston score}{A short, low-dose, non-contrast scan performed before the CCTA, measuring
calcified plaque in the coronaries. It usually adds a second acquisition, which is one reason a whole cardiac
examination's DLP can be higher than a single scan.}

\Tg{ECG gating}{cardiac synchronisation}{The heart moves, so images must be taken at the same phase of the cardiac
cycle. The ECG signal tells the scanner when. Two families:}
\begin{itemize}
\item \textbf{Prospective triggering (step-and-shoot):} radiation is on only during a chosen window (usually
diastole, $\sim$70--75\% of the R--R interval). Low dose. Needs a low, regular heart rate.
\item \textbf{Retrospective gating (helical):} radiation is on throughout the whole cardiac cycle; the phase is
chosen afterwards. Any phase can be reconstructed, and functional analysis is possible --- at 2--5$\times$ the dose.
\end{itemize}
\note{Our 50 exams have a dose level typical of retrospective gating. Switching to prospective triggering, where
the heart rate allows, usually saves more than any kV or mAs change.}

\Tg{Heart rate control}{beta-blockers, nitrates}{Beta-blockers slow the heart (target often $<$60--65 bpm) so
prospective triggering and a short acquisition window become possible; sublingual nitrate dilates the coronaries
for better visibility. Both are dose-relevant even though they are drugs, not scanner settings.}

\T{Scan range / z-coverage}{The length of the patient along the table axis ($z$) that is irradiated. For CCTA
this is typically 12--16 cm per acquisition (from just below the carina to below the heart apex). Dose (DLP) is
directly proportional to it, so trimming the range is free dose reduction.}

% =====================================================================
\section{2 \ \ The scanner settings}

\Tg{kVp}{peak kilovoltage, tube voltage}{The energy of the X-ray beam --- how "hard" or penetrating the X-rays are.
Typical CT values: 70, 80, 100, 120, 140 kV. Dose rises steeply with kV ($\propto \mathrm{kV}^{2.5-3}$), while
iodine contrast \textit{falls} as kV rises. This is the single most powerful lever in contrast-enhanced CT.}

\Tg{mA}{tube current}{How many X-rays per second (the beam's "quantity"). More mA means less image noise and more
dose --- both strictly proportional.}

\Tg{mAs}{tube current $\times$ rotation time}{The total X-ray output per rotation. Dose is exactly linear in mAs.
\textbf{Effective mAs} $=$ mAs / pitch, which is what many consoles display. In our data, mAs varies from 375 to
1345 because of automatic modulation (below).}

\Tg{ATCM}{automatic tube current modulation, e.g. CARE Dose4D, Auto mA}{The scanner varies mA in real time
according to patient size and angle of view, aiming at a constant image quality. This is why mAs correlates with
BMI in our dataset: the machine, not the operator, chose it.}

\T{Pitch}{Table travel per rotation divided by beam width. Pitch $>1$ means the table moves faster than the beam
is wide (less dose, faster scan); pitch $<1$ means overlap (more dose, used in retrospective gating).}

\Tg{Collimation}{beam width}{The total detector width used, e.g. $64 \times 0.6$~mm $=$ 38.4~mm. It affects dose
efficiency and the amount of \textbf{overranging} (extra rotations at the start and end of a helical scan that
irradiate tissue without producing images).}

\T{Rotation time}{Time for one full gantry rotation (0.25--0.5~s). Shorter rotation means better temporal
resolution --- essential for a moving heart.}

\Tg{FOV}{field of view}{The diameter of the reconstructed image. Not a dose parameter in itself, but it is a common
column in scanner exports --- and it is our leading suspect for the ambiguous "SCAN SIZE CM" field in the data.}

\Tg{Reconstruction}{FBP, IR, DLR}{The algorithm turning raw projections into images. Filtered back projection (FBP)
is the classic; iterative reconstruction (IR, e.g. ASIR, SAFIRE, ADMIRE) removes noise and allows 30--50\% dose
reduction at equal noise; deep-learning reconstruction (DLR) goes further. The \textbf{kernel} controls sharpness
versus noise.}

% =====================================================================
\section{3 \ \ Dose quantities --- the units}

\begin{warnbox}
\textbf{The key distinction.} \ctdi{} and DLP describe what the \textit{machine emitted} into a standard plastic
cylinder. They are not the dose absorbed by \textit{this patient}, and not a risk. Effective dose (mSv) is the
risk-weighted estimate; SSDE is the size-corrected patient dose.
\end{warnbox}

\T{Gray (Gy), milligray (mGy)}{Absorbed energy per unit mass (1 Gy = 1 J/kg). The unit of \ctdi.}

\Tg{CTDI$_{100}$, CTDI$_w$}{computed tomography dose index}{Measured in a standard acrylic phantom (16 cm diameter
for head, 32 cm for body) with a 100 mm pencil ionisation chamber. CTDI$_w$ weights centre and periphery
($\tfrac{1}{3}$ centre $+$ \ $\tfrac{2}{3}$ periphery).}

\Tg{\ctdi{} \ (mGy)}{volume CT dose index}{$\mathrm{CTDI}_w/\mathrm{pitch}$ --- the average dose \textbf{per slice}
for a given technique. It measures dose \textit{intensity} and does not depend on how long the scan is. It is the
first of our two prediction targets. In our data: 21--105 mGy, median 47.5.}

\Tg{DLP \ (\mgycm)}{dose--length product}{$\mathrm{DLP} = \ctdi \times$ irradiated length. It measures the
\textbf{total} radiation delivered by the examination, and is what national reference levels are usually expressed
in. Our second prediction target. In our data: 482--2393 \mgycm, median 1080.}

\T{Effective dose \ (mSv)}{A whole-body risk-equivalent number, obtained as $E \approx k \times \mathrm{DLP}$.
For chest CT $k = 0.014$~mSv/\mgycm; for \textbf{cardiac} CT the modern cardiac-specific factor is
$k \approx 0.026$ (0.020--0.035) --- the chest factor underestimates cardiac effective dose by $\sim$46\%.}
\begin{keybox}
\textbf{Our dataset in context.} Median DLP 1080 \mgycm{} $\times$ 0.026 $\approx$ \textbf{28 mSv}.
For comparison: modern prospective CCTA 1--5 mSv; natural background $\approx$2.4 mSv/year; a chest X-ray
$\approx$0.02 mSv; a return transatlantic flight $\approx$0.08 mSv.
\end{keybox}

\Tg{SSDE \ (mGy)}{size-specific dose estimate}{\ctdi{} corrected for the actual patient size, using effective
diameter or water-equivalent diameter (AAPM Reports 204 and 220). A small patient absorbs more dose than the
phantom for the same \ctdi; a large patient absorbs less. Requires a patient dimension we do not currently have.}

\Tg{DRL}{diagnostic reference level}{Not a limit, but a benchmark: the 75th percentile of the dose distribution for
a standard-sized patient across many centres. If your median routinely exceeds the DRL, the protocol should be
reviewed (ICRP 135, EU Directive 2013/59/Euratom). CCTA references: 400 \mgycm{} (PROTECTION VI), UK 170
prospective / 380 retrospective, France 393 / 870.}

\T{Achievable dose}{The \textit{median} of the same distribution --- a more ambitious target than the DRL. The
"peer benchmark" in our app is exactly this idea, computed on your own data per BMI.}

\Tg{ALARA / ALADA}{as low as reasonably achievable / as low as diagnostically acceptable}{The governing principle:
dose should be as low as possible \textit{while still answering the clinical question}. ALADA is the modern
phrasing, and it is precisely the constraint our optimiser implements.}

\T{Stochastic vs deterministic effects}{Stochastic: probability of cancer rises with dose, with no threshold ---
the concern in CT. Deterministic (skin erythema, cataract): only above high thresholds, not reached in diagnostic
CT.}

% =====================================================================
\section{4 \ \ Image quality --- the other half}

\Tg{HU}{Hounsfield unit}{The CT density scale: water $=$ 0, air $=-1000$, dense bone $\approx +1000$. Fat
$\approx -100$, muscle $\approx +50$, contrast-filled coronary lumen $\approx +300$ to $+500$ HU.}

\T{Noise}{The standard deviation of HU in a uniform region (e.g. inside the aorta). Quantum noise falls as
$1/\sqrt{\text{dose}}$: \textbf{to halve the noise you need 4$\times$ the dose}.}

\Tg{SNR / CNR}{signal- and contrast-to-noise ratio}{$\mathrm{CNR} = (\mathrm{HU}_{\text{vessel}} -
\mathrm{HU}_{\text{muscle}})/\text{noise}$. CNR, not dose, is what determines whether a stenosis can be seen. It is
the quantity the phantom loop measures and the quantity the recommendation engine holds constant.}

\T{Spatial resolution}{The smallest detail distinguishable (about 0.3--0.5 mm in CT). Needed to see a small
coronary branch.}

\T{Temporal resolution}{The "shutter speed" for a moving object --- typically half the rotation time
(75--175 ms). Insufficient temporal resolution at high heart rate produces motion artefact.}

\T{Artefacts}{\textbf{Motion} (blurred vessel, from heart rate or arrhythmia); \textbf{blooming} (calcified
plaque or stents look larger than they are); \textbf{beam hardening / streak} (from dense contrast or metal);
\textbf{stair-step} (misregistration between cardiac cycles); \textbf{photon starvation} (excessive noise where
X-rays cannot penetrate --- the risk of using too low a kV in a large patient).}

\T{Diagnostic acceptability score}{A reader's judgement, typically on a 1--4 scale: 4 excellent, 3 diagnostic,
2 limited, 1 non-diagnostic. This is the label the clinical-learning loop uses --- the only ground truth for
"was the dose enough?".}

% =====================================================================
\section{5 \ \ The patient}

\Tg{BMI}{body mass index}{$\mathrm{weight}/\mathrm{height}^2$ in kg/m$^2$. Under 18.5 underweight, 18.5--25 normal,
25--30 overweight, 30--40 obese, $>$40 severely obese. A crude but universally recorded proxy for how much the
patient attenuates X-rays. It is fixed --- hence the project's framing.}

\T{Attenuation}{How much the body absorbs and scatters the beam. Larger and denser bodies need more photons for
the same noise, so dose rises steeply with size.}

\T{Effective diameter / water-equivalent diameter}{Better size descriptors than BMI, measured from the images
themselves; required for SSDE. Worth exporting if it is available.}

\T{Body habitus}{The clinical term for build/shape. Two patients with identical BMI can attenuate differently
(muscle versus fat, broad versus deep chest).}

% =====================================================================
\section{6 \ \ The relationships that drive everything}

\begin{eqbox}
$\ctdi = k \cdot \mathrm{mAs} \cdot (\mathrm{kVp}/120)^{n}$, \ $n \approx 2.5$--$3$\\[3pt]
$\mathrm{DLP} = \ctdi \times \text{irradiated length}$\\[3pt]
$E \ (\mathrm{mSv}) \approx k_{\text{cardiac}} \times \mathrm{DLP}$, \ $k \approx 0.026$\\[3pt]
$\text{noise} \propto 1/\sqrt{\ctdi}$ \qquad $C_{\text{iodine}} \propto \mathrm{kVp}^{-\gamma}$, \ $\gamma \approx 1.35$\\[3pt]
$\mathrm{CNR} \propto C_{\text{iodine}}\sqrt{\ctdi}$
\end{eqbox}

\subsection{Rules of thumb}
\begin{itemize}
\item \textbf{Double the mAs} $\rightarrow$ double the dose, noise $\times 0.71$.
\item \textbf{Halve the noise} $\rightarrow$ 4$\times$ the dose.
\item \textbf{120 $\rightarrow$ 100 kV} at the same mAs $\rightarrow$ dose $\times 0.61$, iodine contrast
$\times 1.3$. At equal iodine CNR, $\sim$40\% less dose.
\item \textbf{120 $\rightarrow$ 80 kV} (slim patients only) $\rightarrow$ $\sim$65\% less dose at equal iodine CNR.
\item \textbf{1 cm less scan length} $\rightarrow$ DLP falls by $1/\text{length}$ ($\sim$3\% for a 30 cm scan);
\ctdi{} does not change.
\item \textbf{Prospective instead of retrospective gating} $\rightarrow$ typically 50--80\% less dose, if the heart
rate allows.
\end{itemize}

% =====================================================================
\section{7 \ \ Greek terms you will meet in the data}

\begin{tabularx}{\columnwidth}{@{}L L@{}}
\toprule
\hdr \textbf{Greek} & \textbf{English}\\
\midrule
ΥΤ (Υπολογιστική Τομογραφία) & CT (computed tomography)\\
ΥΤ ΣΤΕΦΑΝΙΟΓΡΑΦΙΑ & coronary CT angiography (CCTA)\\
Στεφανιαία αγγεία & coronary arteries\\
Σκιαγραφικό & contrast agent\\
Δόση ακτινοβολίας & radiation dose\\
Ακτινοπροστασία & radiation protection\\
Ακτινολόγος / Ακτινοφυσικός & radiologist / medical physicist\\
Ακτινολόγος τεχνολόγος & radiographer, CT technologist\\
Μήκος σάρωσης & scan length\\
Ποιότητα εικόνας & image quality\\
Θόρυβος & (image) noise\\
Ομοίωμα & phantom\\
Επίπεδα αναφοράς (DRL) & diagnostic reference levels\\
\bottomrule
\end{tabularx}

% =====================================================================
\section{8 \ \ How each term shows up in the app}

\begin{tabularx}{\columnwidth}{@{}l L@{}}
\toprule
\hdr \textbf{Term} & \textbf{Where it appears}\\
\midrule
BMI & input, fixed; also defines the minimum kV rule and the peer benchmark\\
kVp & input and the main optimisation lever\\
mAs & input; the engine computes the compensating mAs for a kV change\\
Scan length & input; scales DLP proportionally\\
\ctdi{} & prediction target 1, with a 90\% interval\\
DLP & prediction target 2, compared with your reference level\\
CNR & the constraint: "keep CNR at $\ge$ x\% of the current plan"\\
$\gamma$ (iodine exponent) & literature 1.35, or calibrated by your phantom campaign\\
Phantom & the Bayesian-optimisation tab: propose scan $\rightarrow$ measure CNR $\rightarrow$ repeat\\
Diagnostic score 1--4 & the clinical-learning tab: the reward the bandit learns from\\
DRL & the reference line on the dose charts (default 400 \mgycm)\\
\bottomrule
\end{tabularx}

% =====================================================================
\section{10 \ \ Anatomy \& findings you will hear}

\Tg{LAD / LCx / RCA}{left anterior descending, left circumflex, right coronary artery}{The three main coronary
arteries. The LAD supplies the largest territory, so a proximal LAD stenosis is the most feared finding. The left
main (LM) is the short trunk that splits into LAD and LCx.}

\T{Stenosis}{A narrowing of the artery, reported in percent of the lumen diameter ($<$50\% non-obstructive,
$\ge$50\% obstructive, $\ge$70\% significant). Seeing a 50\% versus a 70\% narrowing is exactly the task that image
quality has to support.}

\T{Plaque}{Deposit in the artery wall. \textbf{Calcified} plaque is bright and causes blooming artefact;
\textbf{non-calcified (soft)} plaque is low-density and needs good contrast-to-noise to detect --- another reason
noise matters more than it might seem.}

\T{Stent}{A metal mesh keeping an artery open. Metal creates artefacts, so stented segments usually need a
sharper kernel and lower noise (i.e. more dose).}

\Tg{CAD-RADS}{coronary artery disease reporting and data system}{The standardised 0--5 reporting scale for CCTA,
plus a modifier for non-diagnostic studies (N). "CAD-RADS N" is, in effect, a wasted dose.}

\Tg{FFR-CT}{CT-derived fractional flow reserve}{A computational-fluid-dynamics analysis of the CCTA images
estimating whether a stenosis actually restricts flow. It demands good image quality, so it raises the cost of
being too aggressive on dose.}

% =====================================================================
\section{11 \ \ Data \& IT terms}

\Tg{DICOM}{digital imaging and communications in medicine}{The file and network standard for medical images. Each
image carries a header of tags --- including KVP (0018,0060), Exposure (0018,1152), CTDIvol (0018,9345) and the
scan parameters we need.}

\Tg{RDSR}{radiation dose structured report}{A DICOM object the scanner produces per examination, containing the
per-series \ctdi, DLP, kV, mAs and scan length. \textbf{This is the ideal data source for the app} --- far better
than typing numbers off the console, and it also contains the true scan length.}

\Tg{PACS / RIS}{picture archiving and communication system / radiology information system}{Where images and reports
live. Dose-management platforms (e.g. DoseWatch, Radimetrics, OpenREM) harvest RDSR from PACS automatically.}

\T{Pseudonymisation}{Replacing patient identifiers with a code. Our clinical-learning log stores only a
pseudonymised ID, BMI and dose values, and never leaves the local computer --- keep it that way for GDPR reasons.}

% =====================================================================
\section{9 \ \ Who does what}

\T{Radiologist}{The doctor who reads the images and answers the clinical question. Owns the judgement of whether
image quality is diagnostic --- the score the bandit learns from.}

\T{Medical physicist}{Responsible for dose measurement, phantom work, scanner QA and protocol optimisation.
Your main partner for the phantom campaign, and the person who signs off any protocol change.}

\T{Radiographer / CT technologist}{Runs the scan, positions the patient, selects the protocol on the console.
The person who would use the "Predict" and "Reduce dose" tabs day to day.}

\T{Cardiologist}{Usually the referrer, and the reader in some centres.}

\begin{warnbox}
\textbf{Two sentences worth memorising.}\\
1. \ctdi{} is dose \textit{intensity} (per slice); DLP is \textit{total} dose (intensity $\times$ length); mSv is
\textit{risk}.\\
2. Dose without image quality is meaningless in both directions --- more dose is not "safer imaging", and less dose
is not an improvement if the study has to be repeated.
\end{warnbox}

\vfill
\note{Sources: Stocker et al., PROTECTION VI, Eur Heart J 2018 (kV by BMI, 400 \mgycm{} reference);
Trattner et al., \textit{Cardiac-specific conversion factors}, JACC Cardiovasc Imaging 2018 ($k = 0.026$);
AAPM Reports 204 \& 220 (SSDE); ICRP 103/135; EU Directive 2013/59/Euratom. Typical values are
indicative --- your scanner, protocol and phantom are the authority.}

\end{multicols}
\end{document}
